% @@STATS_BEGIN@@ — auto-generated by habor-analyze, do not edit manually
% Auto-generated by habor-analyze — do not edit manually.
% Re-generate: uv run habor-analyze studies

% ── General ──
\providecommand{\NumIncludedBenchmarks}{54}
\providecommand{\NumHaborMixTasks}{160}
\providecommand{\NumModels}{8}
\providecommand{\NumAgents}{4}

% ── Within-family: model vs agent effect ──
\providecommand{\AnthropicModelEffectMean}{0.236}
\providecommand{\AnthropicAgentEffectMean}{0.093}
\providecommand{\AnthropicModelAgentRatio}{2.5}
\providecommand{\AnthropicModelWins}{48}
\providecommand{\AnthropicTotalBenchmarks}{54}
\providecommand{\AnthropicModelWinPct}{89}
\providecommand{\GoogleModelEffectMean}{0.084}
\providecommand{\GoogleAgentEffectMean}{0.111}
\providecommand{\GoogleModelAgentRatio}{0.8}
\providecommand{\GoogleModelWins}{31}
\providecommand{\GoogleTotalBenchmarks}{54}
\providecommand{\GoogleModelWinPct}{57}
\providecommand{\OpenAIModelEffectMean}{0.267}
\providecommand{\OpenAIAgentEffectMean}{0.164}
\providecommand{\OpenAIModelAgentRatio}{1.6}
\providecommand{\OpenAIModelWins}{42}
\providecommand{\OpenAITotalBenchmarks}{54}
\providecommand{\OpenAIModelWinPct}{78}

% ── Adjusted effects ──
\providecommand{\ModelEffectSpan}{2.80}
\providecommand{\AgentEffectSpan}{0.17}
\providecommand{\TopModel}{gpt-5.4}
\providecommand{\TopModelEffect}{0.52}
\providecommand{\BottomModel}{gpt-5-nano}
\providecommand{\BottomModelEffect}{-2.28}
\providecommand{\TopAgent}{claude-code}
\providecommand{\TopAgentEffect}{-0.21}
\providecommand{\BottomAgent}{gemini-cli}
\providecommand{\BottomAgentEffect}{-0.39}

% ── Agent lift vs terminus-2 ──
\providecommand{\claudecodeMeanDelta}{0.263}
\providecommand{\claudecodeWinRate}{64}
\providecommand{\codexMeanDelta}{-0.005}
\providecommand{\codexWinRate}{70}
\providecommand{\geminicliMeanDelta}{-0.087}
\providecommand{\geminicliWinRate}{56}

% ── Terminus delta by model×agent (top pairs) ──
\providecommand{\TerminusgptFiveFourcodexDelta}{0.812}
\providecommand{\TerminusgptFiveFourcodexWinRate}{94}
\providecommand{\TerminusclaudehaikuFourFiveTwoZeroTwoFiveOneZeroZeroOneclaudecodeDelta}{0.684}
\providecommand{\TerminusclaudehaikuFourFiveTwoZeroTwoFiveOneZeroZeroOneclaudecodeWinRate}{81}
\providecommand{\TerminusgptFiveminicodexDelta}{0.543}
\providecommand{\TerminusgptFiveminicodexWinRate}{74}
\providecommand{\TerminusclaudesonnetFourSixclaudecodeDelta}{0.091}
\providecommand{\TerminusclaudesonnetFourSixclaudecodeWinRate}{65}
\providecommand{\TerminusgeminiThreeOnepropreviewgeminicliDelta}{0.039}
\providecommand{\TerminusgeminiThreeOnepropreviewgeminicliWinRate}{57}
\providecommand{\TerminusclaudeopusFourSixclaudecodeDelta}{0.013}
\providecommand{\TerminusclaudeopusFourSixclaudecodeWinRate}{44}
\providecommand{\TerminusgeminiThreeflashpreviewgeminicliDelta}{-0.212}
\providecommand{\TerminusgeminiThreeflashpreviewgeminicliWinRate}{56}
\providecommand{\TerminusgptFivenanocodexDelta}{-1.371}
\providecommand{\TerminusgptFivenanocodexWinRate}{43}

% ── Agent lift table (paper, short names) ──
\providecommand{\GPTFiveFourBase}{.510}
\providecommand{\GPTFiveFourDelta}{$+$.184}
\providecommand{\GPTFiveFourLift}{94}
\providecommand{\ClaudeHaikuBase}{.390}
\providecommand{\ClaudeHaikuDelta}{$+$.107}
\providecommand{\ClaudeHaikuLift}{81}
\providecommand{\GPTMiniBase}{.405}
\providecommand{\GPTMiniDelta}{$+$.121}
\providecommand{\GPTMiniLift}{74}
\providecommand{\ClaudeSonnetBase}{.568}
\providecommand{\ClaudeSonnetDelta}{$+$.036}
\providecommand{\ClaudeSonnetLift}{65}
\providecommand{\GeminiProBase}{.642}
\providecommand{\GeminiProDelta}{$+$.036}
\providecommand{\GeminiProLift}{57}
\providecommand{\ClaudeOpusBase}{.612}
\providecommand{\ClaudeOpusDelta}{$+$.025}
\providecommand{\ClaudeOpusLift}{44}
\providecommand{\GeminiFlashBase}{.567}
\providecommand{\GeminiFlashDelta}{$-$.006}
\providecommand{\GeminiFlashLift}{56}
\providecommand{\GPTNanoBase}{.261}
\providecommand{\GPTNanoDelta}{$-$.015}
\providecommand{\GPTNanoLift}{43}

% ── Benchmark headroom ──
\providecommand{\NumSaturatedBenchmarks}{11}
\providecommand{\NumHardBenchmarks}{7}
\providecommand{\MeanBestScore}{0.739}

% ── Per-domain headroom stats ──
\providecommand{\NumAgentsToolsAndSystemsSaturated}{0}
\providecommand{\NumAgentsToolsAndSystemsHard}{1}
\providecommand{\NumAgentsToolsAndSystemsBenchmarks}{9}
\providecommand{\MeanAgentsToolsAndSystemsBestScore}{67}
\providecommand{\NumDataAndAnalyticsSaturated}{0}
\providecommand{\NumDataAndAnalyticsHard}{1}
\providecommand{\NumDataAndAnalyticsBenchmarks}{2}
\providecommand{\MeanDataAndAnalyticsBestScore}{51}
\providecommand{\NumKnowledgeAndLongContextSaturated}{2}
\providecommand{\NumKnowledgeAndLongContextHard}{0}
\providecommand{\NumKnowledgeAndLongContextBenchmarks}{4}
\providecommand{\MeanKnowledgeAndLongContextBestScore}{83}
\providecommand{\NumMathematicsAndReasoningSaturated}{3}
\providecommand{\NumMathematicsAndReasoningHard}{0}
\providecommand{\NumMathematicsAndReasoningBenchmarks}{6}
\providecommand{\MeanMathematicsAndReasoningBestScore}{92}
\providecommand{\NumMultimodalSaturated}{0}
\providecommand{\NumMultimodalHard}{0}
\providecommand{\NumMultimodalBenchmarks}{1}
\providecommand{\MeanMultimodalBestScore}{70}
\providecommand{\NumProfessionalDomainsSaturated}{0}
\providecommand{\NumProfessionalDomainsHard}{0}
\providecommand{\NumProfessionalDomainsBenchmarks}{5}
\providecommand{\MeanProfessionalDomainsBestScore}{75}
\providecommand{\NumSafetyAndSecuritySaturated}{1}
\providecommand{\NumSafetyAndSecurityHard}{0}
\providecommand{\NumSafetyAndSecurityBenchmarks}{1}
\providecommand{\MeanSafetyAndSecurityBestScore}{97}
\providecommand{\NumScientificResearchSaturated}{0}
\providecommand{\NumScientificResearchHard}{1}
\providecommand{\NumScientificResearchBenchmarks}{8}
\providecommand{\MeanScientificResearchBestScore}{67}
\providecommand{\NumSoftwareEngineeringSaturated}{5}
\providecommand{\NumSoftwareEngineeringHard}{4}
\providecommand{\NumSoftwareEngineeringBenchmarks}{18}
\providecommand{\MeanSoftwareEngineeringBestScore}{74}
\providecommand{\MostHeadroomDomainOne}{Data \& Analytics}
\providecommand{\MostHeadroomDomainOneScore}{51}
\providecommand{\MostHeadroomDomainTwo}{Agents, Tools \& Systems}
\providecommand{\MostHeadroomDomainTwoScore}{67}
\providecommand{\MostHeadroomDomainThree}{Scientific Research}
\providecommand{\MostHeadroomDomainThreeScore}{67}
% ── SE hardest benchmarks (sorted by score) ──
\providecommand{\SEHardBenchmarkOne}{AlgoTune}
\providecommand{\SEHardBenchmarkOneScore}{37}
\providecommand{\SEHardBenchmarkTwo}{FeatureBench}
\providecommand{\SEHardBenchmarkTwoScore}{41}
\providecommand{\SEHardBenchmarkThree}{BigCodeBench}
\providecommand{\SEHardBenchmarkThreeScore}{45}
\providecommand{\SEHardBenchmarkFour}{GSO}
\providecommand{\SEHardBenchmarkFourScore}{45}

% ── Effective dimensionality ──
\providecommand{\ParticipationRatio}{1.3}
\providecommand{\ComponentsForNinety}{2}
\providecommand{\ComponentsForNinetyFive}{3}
\providecommand{\TopOneVariance}{89}
\providecommand{\TopThreeVariance}{97}

% ── Domain correlation ──
\providecommand{\WithinDomainMedianCorr}{0.68}
\providecommand{\CrossDomainMedianCorr}{0.62}
\providecommand{\WithinDomainFracAboveSeven}{44}
\providecommand{\CrossDomainFracAboveSeven}{38}
\providecommand{\WithinAgentsToolsAndSystemsMedianCorr}{0.65}
\providecommand{\WithinKnowledgeAndLongContextMedianCorr}{0.58}
\providecommand{\WithinMathematicsAndReasoningMedianCorr}{0.59}
\providecommand{\WithinProfessionalDomainsMedianCorr}{0.71}
\providecommand{\WithinScientificResearchMedianCorr}{0.48}
\providecommand{\WithinSoftwareEngineeringMedianCorr}{0.69}

% ── Greedy benchmark selection ──
\providecommand{\GreedyBenchmarksBelowSeven}{13}
\providecommand{\GreedyFirstAboveSeven}{gso}
\providecommand{\GreedyFirstAboveSevenStep}{14}

% ── HaborMix difficulty composition ──
\providecommand{\HaborMixMediumTasks}{99}
\providecommand{\HaborMixEasyTasks}{44}
\providecommand{\HaborMixHardTasks}{16}
\providecommand{\HaborMixFrontierTasks}{1}

% ── AIME case study ──
\providecommand{\AimeGptFourCodexBase}{0.711}
\providecommand{\AimeGptFourCodexAgent}{0.989}
\providecommand{\AimeGptFourCodexDeltaPP}{27.8}
\providecommand{\AimeGptMiniCodexBase}{0.889}
\providecommand{\AimeGptMiniCodexAgent}{0.967}
\providecommand{\AimeGptMiniCodexDeltaPP}{7.8}
\providecommand{\AimeGptNanoCodexBase}{0.756}
\providecommand{\AimeGptNanoCodexAgent}{0.739}
\providecommand{\AimeGptNanoCodexDeltaPP}{-1.7}

% ── Figure paths (all under paper/figs/main/quantitative/) ──
\providecommand{\FigHeadroomByDomain}{../figs/main/quantitative/bench\_progress\_and\_headroom}

% @@STATS_END@@

\subsubsection{Benchmarks}

%%% !todo: need to make sure all the figures and data being used is from the finalized benchmarks.

%%% for each part of these, it can be included in the appendix for all the math equations, data details, and more figures.

%%%%%%%%%%%%%%%%%%%%%%%%%%%%%%%%%%%%%%%%%%%%%%%%%%%%%%%%%%%%%%%%%%%%%%%%%%%%%%%
\paragraph{How fast are models improving over time?}
% TODO: progress-over-time figure (needs model release dates in data).
% \lin{We need one main figure here: progress by the years of some main
% benchmarks that we care and can make apple-to-apple comparison. Each domain
% should share the same color - different benchmarks should have different
% color intensity.}

%%%%%%%%%%%%%%%%%%%%%%%%%%%%%%%%%%%%%%%%%%%%%%%%%%%%%%%%%%%%%%%%%%%%%%%%%%%%%%%
\paragraph{Which benchmarks still challenge frontier systems?}
Current evaluations exhibit a severe saturation problem. Among the
\NumIncludedBenchmarks{} benchmarks analyzed (\outlinecomment{xiangning: these need to use the finalized benchmarks}),
about one-fifth (\NumSaturatedBenchmarks{}) have already been effectively
solved (best model with agents has scores $>$95\%). This saturation spans traditionally
``difficult'' domains, including Mathematics \& Reasoning (e.g., AIME, KUMO)
and Knowledge \& Long Context (e.g., GPQA Diamond). Consequently, these
benchmarks now offer minimal headroom to differentiate the next generation
of frontier models.

The remaining challenge lies unevenly across domains. While areas like
\MostHeadroomDomainOne{} and \MostHeadroomDomainTwo{} retain substantial
headroom, Software Engineering emerges as the most structurally heterogeneous
domain. It simultaneously contains the highest concentration of saturated
tasks (e.g., HumanEvalFix, CompileBench) and the highest concentration of
unsolved challenges (e.g., \SEHardBenchmarkOne{}, \SEHardBenchmarkTwo{},
where state-of-the-art systems still score below 50\%). This dichotomy
highlights a fundamental shift in AI capabilities: while isolated,
well-specified coding tasks (like bug fixing or snippet generation) are
largely solved, open-ended, system-level software engineering remains a
true frontier.

\begin{figure}[t]
  \centering
  \includegraphics[width=\linewidth]{\FigHeadroomByDomain}
  \caption{Best system score per benchmark, colored by domain. The shaded area
    represents headroom---room for improvement. Benchmarks are sorted within
    each domain by best score.}
  \label{fig:headroom-by-domain}
\end{figure}

%%%%%%%%%%%%%%%%%%%%%%%%%%%%%%%%%%%%%%%%%%%%%%%%%%%%%%%%%%%%%%%%%%%%%%%%%%%%%%%
\paragraph{Are benchmark scores predictable, and does agent scaffolding alter this structure?}

Recent work shows model-only benchmark scores are highly
redundant and predictable via low-rank matrix
completion~\cite{papailiopoulos2026benchpress}. We find a
similar low-rank structure persists in our setting, but it
fundamentally changes when introducing advanced agent
scaffolds. Under the minimal Terminus-2 baseline, a single
principal component explains \AppSVDModelOnlyPCOne{}\% of
variance (comparable to the 71\% reported by BenchPress on
scaffold-free evaluations). When advanced scaffolds are
included, a second orthogonal component emerges
(\AppSVDFullPCTwo{}\% variance) that perfectly separates
advanced-scaffold systems from baseline systems.
This indicates agent scaffolding is not a uniform score
multiplier, but introduces a structurally distinct capability
dimension (Appendix~\ref{app:benchpress} details SVD
methodology and task-level predictability limitations).

Despite this added complexity, the combined model-agent matrix
remains highly predictable. Using a matrix-completion approach
(Appendix~\ref{app:benchpress}), we find that many established
benchmarks are largely redundant: scores on
\BenchPressRedundantOne{}
(\BenchPressRedundantOneScore{}\% median absolute percentage
error), \BenchPressRedundantTwo{}
(\BenchPressRedundantTwoScore{}\%), and
\BenchPressRedundantThree{} (\BenchPressRedundantThreeScore{}\%)
can be reliably reconstructed from a sparse subset of other
evaluations. Conversely, benchmarks that resist
prediction---such as \BenchPressUniqueOne{}
(\BenchPressUniqueOneScore{}\% error), \BenchPressUniqueTwo{}
(\BenchPressUniqueTwoScore{}\%), and \BenchPressUniqueThree{}
(\BenchPressUniqueThreeScore{}\%)---measure capabilities orthogonal to the
dominant axes. Notably, AIME---despite being considered a difficult
mathematics benchmark---is among the most predictable
(\BenchPressRedundantOneScore{}\%-class MedAPE), confirming that
difficulty and uniqueness are decoupled: to meaningfully differentiate
future agentic systems, evaluation design must pivot from creating
harder static questions toward tasks that demand open-ended
exploration and long-horizon planning.

%%%%%%%%%%%%%%%%%%%%%%%%%%%%%%%%%%%%%%%%%%%%%%%%%%%%%%%%%%%%%%%%%%%%%%%%%%%%%%%
\subsubsection{Agent Harness}
\label{sec:agent-harness}
\outlinecomment{0.5 pages}

\paragraph{Does the model or the agent matter more?}
When deploying an agentic system, practitioners must choose both a base
model and an agent scaffold. To quantify their relative importance, we
compare the variance introduced by swapping models versus swapping agents
within the same provider family (Appendix~\ref{app:within-family}).

The results strongly favor model selection. In both the Anthropic family
(model-dominant on \AnthropicModelWinPct\% of benchmarks, ratio
\AnthropicModelAgentRatio{}$\times$) and the OpenAI family
(\OpenAIModelWinPct\%, \OpenAIModelAgentRatio{}$\times$), upgrading the
underlying model yields a substantially larger score change than switching
the agent harness. The Google family is the exception: with only two
models and a ratio of \GoogleModelAgentRatio{}$\times$, agent and model
effects are comparable (\GoogleModelWinPct\% model-dominant). When
adjusting for baseline difficulty across all \NumModels{} models and
\NumAgents{} agents (Appendix~\ref{app:adjusted-effects}), the variance
attributed to the base model exceeds that of the agent scaffold by over
an order of magnitude. Thus, while sophisticated agent harnesses garner
significant attention, the raw capability of the underlying foundation
model remains the primary bottleneck for system-level performance.

%%%%%%%%%%%%%%%%%%%%%%%%%%%%%%%%%%%%%%%%%%%%%%%%%%%%%%%%%%%%%%%%%%%%%%%%%%%%%%%
\paragraph{Does the agent harness help, and when?}
We compare each agent scaffold against the Terminus~2 neutral baseline across
all benchmarks. Table~\ref{tab:agent-lift} holds the agent fixed within each
provider family and varies the underlying model from strong to weak.

\begin{table}[t]\small
\centering
\caption{Agent lift over Terminus~2 by model tier within each provider family.
  Agent~$\Delta$ is the mean score change; Positive Lift is the fraction of
  benchmarks where the agent scaffold strictly improves upon the neutral
  harness score.}
\label{tab:agent-lift}
\begin{tabular*}{\linewidth}{@{\extracolsep{\fill}}llccc@{}}
\toprule
Family & Model (tier) & Base & Agent $\Delta$ & Positive Lift \\
\midrule
\multirow{3}{*}{OpenAI + Codex}
  & GPT-5.4 (strong)  & \GPTFiveFourBase & \GPTFiveFourDelta & \GPTFiveFourLift\% \\
  & GPT-5-Mini (mid)   & \GPTMiniBase & \GPTMiniDelta & \GPTMiniLift\% \\
  & GPT-5-Nano (weak)  & \GPTNanoBase & \GPTNanoDelta & \GPTNanoLift\% \\
\midrule
\multirow{3}{*}{Anthropic + Claude Code}
  & Claude-Opus (strong)  & \ClaudeOpusBase & \ClaudeOpusDelta & \ClaudeOpusLift\% \\
  & Claude-Sonnet (mid)   & \ClaudeSonnetBase & \ClaudeSonnetDelta & \ClaudeSonnetLift\% \\
  & Claude-Haiku (weak)   & \ClaudeHaikuBase & \ClaudeHaikuDelta & \ClaudeHaikuLift\% \\
\midrule
\multirow{2}{*}{Google + Gemini CLI}
  & Gemini Pro (strong) & \GeminiProBase & \GeminiProDelta & \GeminiProLift\% \\
  & Gemini Flash (mid)  & \GeminiFlashBase & \GeminiFlashDelta & \GeminiFlashLift\% \\
\bottomrule
\end{tabular*}
\end{table}

Agent scaffolding does not uniformly improve model
performance. A necessary condition is that the base model
exceeds a \emph{competence floor} set by the orchestration
overhead of the scaffold (planning, tool dispatch, context
management); below this floor, agentic wrapping degrades
results---GPT-5-Nano under Codex improves on only \GPTNanoLift\% of
benchmarks, and Gemini~Flash under Gemini~CLI shows essentially zero
mean delta (\GeminiFlashDelta{}) despite improving on \GeminiFlashLift\%
of benchmarks individually (Table~\ref{tab:agent-lift}).

Above this floor, the magnitude and distribution of gains
depend on the scaffold--model pairing in a structured way.
We observe two contrasting regimes.
\textbf{Capability amplification}, exemplified by Codex,
where gains scale \emph{with} base-model strength
(Table~\ref{tab:agent-lift}), consistent with a scaffold
that exposes capabilities the model already possesses but
cannot deploy in single-turn evaluation.
\textbf{Gap closing}, exemplified by Claude~Code, where
gains scale \emph{inversely} with base-model strength,
consistent with a scaffold that compensates for weaknesses
through retry and decomposition strategies.
Gemini~CLI falls between the two, providing a modest,
roughly uniform lift across tiers.

